\documentclass[a4paper,12pt]{article}

\usepackage[utf8]{inputenc}
\usepackage[T1]{fontenc}
\usepackage[a4paper, margin=2.5cm]{geometry}
\usepackage{hyperref}

\title{Report}
\author{Jakub Bláha}
\date{\today}

\begin{document}

\maketitle
\tableofcontents
\newpage

\section{Introduction}

This section defines the goals of the work and the motivation behind it as well as the general steps taken towards achieving the goal.

\section{Formalism requirements}

This section extends on the motivation and concretely defienes what the features of the formalism need to be in order to express everything that is needed to be expressed.

\section{Formalism design process}

This section lists the steps taken towards achieving the final version of the formalism.

\section{Known limitations and possible extensions}

This section lists limitations -- that is what the formalism is unable to express -- and outlines how these limitations can be addressed in future versions.

\section{Compatibility with existing formal verification tools}

This section is on how the requirements and test cases expressed using the formalism are convertible to known widely used formal bases such as temporal logic and timed automata.

\section{Requirement and test cases conversion examples}

This sections lists a few examples of requirements and how they can be written in the formalism.

\section{Future work}

This section outline the steps that will be next taken in order to implement the wanted tool.

\end{document}
